\documentclass{minimal}
\usepackage{xltxtra}
\usepackage{fontspec}
\usepackage{polyglossia}
\setmainlanguage{russian}
\setotherlanguage{english}
\setmainfont{Times New Roman}
\newfontfamily{\cyrillicfont}{Times New Roman}
\usepackage{amsmath}
\begin{document}
    Эту задачу надо решать исходя из энергии системы в начальный момент и на момент, когда пружина пытается оторвать нижнюю шайбу от стола. Рассмотрим начальный момент, в этом случае пружина сжата весом двух грузов, груза $m$ и верхней
    шайбы $M$, при этом деформация пружины:
    \[
        x_1 = \frac{(M + m)g}{k}
    \]
    где $k$ - коэффициент упругости пружины, а $g$ - ускорение свободного падения. Энергия сжатой пружины:
    \[
        E_1 = \frac{k x_1^2}{2}
    \]
    После того как груз $m$ сброшен, верхняя шайба начинает двигаться вверх, по инерции проходит точку, где она не сжимает пружину и начинает растягивать пружину. Для того, что нижняя шайба начала отрываться от стола, пружина должна
    действовать на нее с силой:
    \[
        F = Mg
    \]
    для этого пружина должна быть растянута на:
    \[
        x_2 = \frac{Mg}{k}
    \]
    Энергия необходимая для растягивания пружины на $x_2$:
    \[
        E_2=\frac{k x_2^2}{2}
    \]
    Кроме энергии пружины у нас изменяется потенциальная энергия верхней шайбы, она поднимется на $x_1 + x_2$ и её изменение потенциальной энергии составит:
    \[
        E_{\text{п}} = Mg(x_1 + x_2)
    \]
    Составим баланс энергии:
    \[
        \frac{k x_1^2}{2} \geq \frac{k x_2^2}{2} + Mg(x_1 + x_2)
    \]
    Перенесем энергию пружины в одну сторону:
    \[
        \frac{k x_1^2}{2} - \frac{k x_2^2}{2} \geq Mg(x_1 + x_2)
    \]
    Вынесем $\frac{k}{2}$ за скобку:
    \[
        \frac{k}{2}(x_1^2 - x_2^2) \geq Mg(x_1 + x_2)
    \]
    Разложим $x_1^2 - x_2^2$ по формуле разности квадратов:
    \[
        \frac{k}{2}(x_1 - x_2)(x_1 + x_2) \geq Mg(x_1 + x_2)
    \]
    Сократим на $(x_1 + x_2)$:
    \[
        \frac{k}{2}(x_1 - x_2) \geq Mg
    \]
    Подставим вместо $x_1$ и $x_2$ их значения:
    \[
        \frac{k}{2}(\frac{(M + m)g}{k} - \frac{Mg}{k}) \geq Mg
    \]
    Умножим на $k$ и 2 перенесем вправо:
    \[
        (M + m)g - Mg \geq 2Mg
    \]
    Разделим обе части на $g$ и приведем подобные:
    \[
        m \geq 2M
    \]
    Откуда
    \[
        M \leq \frac{m}{2}
    \]
    Подставим $m$
    \[
        M \leq \frac{0.5\ \text{кг}}{2}
    \]
    \[
        M \leq 0.25\ \text{кг}
    \]

\end{document}
